% !TEX root = ../main.tex
\section{Introduction}
\label{sec:introduction}

Agents that execute complex, multi-step tasks in tandem with humans are increasingly central to productivity and decision-making. In goal-oriented workflows---e.g., writing assistants where users draft, revise, and publish, or product flows where users sign up, subscribe, or cancel---success is measured by \emph{outcome metrics}: whether the user published the article, subscribed to the service, or completed the intended action, rather than by turn-level correctness alone. A central question for builders and researchers is: How do system changes (e.g., a new or improved tool for one goal, or a prompt change) impact these end-user outcomes? Answering it is essential for reliable A/B testing, safe deployment, and iterative improvement of agent systems.

Existing evaluation approaches face two main limitations when applied to long-horizon, human-in-the-loop agent workflows \citep{kohavi2020trustworthy, takanobu2020goal, imbens2015causal}. First, small, local changes get diluted in long trajectories: a better tool for ``create diagram'' may improve success on that subgoal, but the signal is mixed with many other steps and goals, so aggregate metrics (e.g., overall task completion) are noisy and poorly attributable to the change. Second, goal boundaries are fuzzy and context-dependent: users switch between goals in arbitrary order (e.g., draft abstract, then references, then back to abstract), and segmenting trajectories into goal-level units often requires ad hoc clustering or heuristics, which undermines both reproducibility and interpretability of metrics.

We propose to address these challenges with a synthetic simulator grounded in a nested Markov chain model (see Method, Figure~\ref{fig:nested-markov}). The top-level chain includes a \emph{start} state, \emph{goal} states (e.g., write abstract, create table, make diagram, research references, run script, write conclusion) that can be visited in arbitrary order, outcome states (\emph{publish}, \emph{subscribe}), and terminal states (\emph{finished}, \emph{abandoned}). The nested chain, per goal, represents sub-states \emph{succeeded}, \emph{continue}, and \emph{failed}. Goal-specific tools influence the transition matrix of the nested chain, and the latent outcome of each goal's subprocess feeds back into the top-level chain (e.g., biasing the next goal). This structure allows us to (1) generate long-horizon trajectories with a clear goal-level decomposition, (2) vary system components (e.g., tools or prompts) at the goal level and study their effect on outcome metrics, and (3) avoid ad hoc trajectory segmentation by treating goals as explicit states in the model.

Under a finite set of user personas, the model yields closed-form formulas for outcome probabilities, sample size, and experiment design. We emphasize four consequences. When an outcome (e.g., publish) correlates positively with a rarer monetization event (e.g., subscribe) and occurs more frequently, using it as the A/B metric can \emph{reduce} the sample size required for a given power, enabling faster and cheaper experiments (\emph{proxies}). The best split of a fixed trajectory budget between two systems is \emph{not} 50/50: allocating more to the noisier system minimizes the variance of the estimated difference and maximizes power, a counterintuitive result that speeds up experiments (\emph{optimal A/B allocation}). Calibrating the chain to data yields per-goal and per-persona sensitivities; ranking by impact per unit effort identifies which tools and prompt enhancements have the highest improvement potential, thereby maximizing return on scarce engineering effort and organizational velocity (\emph{leverage}). A perfect labelling algorithm is not needed; EM for transition estimation corrects for mislabelling and reduces variance (\emph{robustness to label noise}; see \S\ref{sec:method:inference}).

\paragraph{Problem statements.}
We formalize the setting as follows.

\begin{enumerate}
  \item \emph{Outcome attribution.} Given trajectories of agent-human interaction and a discrete outcome (e.g., publish, subscribe, cancel), estimate the causal or associative effect of a system intervention (e.g., changing a tool or prompt for a specific goal) on the probability or rate of that outcome, despite long trajectories and confounders from other goals and steps.

  \item \emph{Goal-aware evaluation.} Define a reproducible, goal-level decomposition of trajectories so that metrics (success, continue, fail) can be defined per goal and aggregated in a way that reflects the user's multi-goal, non-sequential workflow, without relying on ad hoc clustering of turns.

  \item \emph{Synthetic, scalable evaluation.} Design a simulator that (a) mimics goal-oriented, human-in-the-loop behavior, (b) exposes system levers (e.g., per-goal tools and transition dynamics), and (c) generates enough data to measure the impact of small system changes on outcome metrics that would otherwise be diluted in real long trajectories.
\end{enumerate}

We address (1)--(3) with a nested Markov chain simulator: under a finite set of user personas, outcome attribution, sample size, and design levers are closed form, and the model can be calibrated to labelled trajectories (with EM correcting for label noise). We focus on this synthetic setting and discrete personas; validation on real production logs is discussed in the outlook. Controlled experiments and synthetic evaluation are standard for measuring treatment effects \citep{kohavi2020trustworthy, rubin2005causal}; we build on this by making goals and outcome attribution explicit in the chain.

\paragraph{Outline.}
We give background (\S\ref{sec:background}), then present the nested Markov chain and simulator (\S\ref{sec:method}) and experiments (\S\ref{sec:experiments}). We then detail design levers: leverage (\S\ref{sec:leverage}), minimal detectable effect (\S\ref{sec:mde}), and optimal A/B allocation (\S\ref{sec:optimal-ab}). We conclude (\S\ref{sec:conclusion}) and discuss limitations and outlook (\S\ref{sec:outlook}).
